\documentclass[a4paper,11pt]{article}

%%% Packages
\usepackage[utf8]{inputenc}
\usepackage[T1]{fontenc}
\usepackage{amsmath,amssymb,amsthm}
\usepackage{mathtools}
\usepackage{stmaryrd}       % \llbracket, \rrbracket
\usepackage{booktabs}
\usepackage{array}
\usepackage{enumitem}
\usepackage{hyperref}
\usepackage[margin=2.5cm]{geometry}
\usepackage{xcolor}

%%% Theorem environments
\newtheorem{theorem}{Theorem}[section]
\newtheorem{lemma}[theorem]{Lemma}
\newtheorem{proposition}[theorem]{Proposition}
\newtheorem{corollary}[theorem]{Corollary}
\newtheorem{claim}[theorem]{Claim}
\newtheorem{conjecture}[theorem]{Conjecture}
\theoremstyle{definition}
\newtheorem{definition}[theorem]{Definition}
\newtheorem{example}[theorem]{Example}
\newtheorem{remark}[theorem]{Remark}

%%% Notation macros (matching companion paper)
\newcommand{\ALC}{\ensuremath{\mathcal{ALC}}}
\newcommand{\ALCI}{\ensuremath{\mathcal{ALCI}}}
\newcommand{\ALCIRCC}[1]{\ensuremath{\mathcal{ALCI}_{\mathit{RCC#1}}}}
\newcommand{\NR}{\ensuremath{N_R}}
\newcommand{\NC}{\ensuremath{N_C}}
\newcommand{\DeltaI}{\ensuremath{\Delta^{\mathcal{I}}}}
\newcommand{\interp}{\ensuremath{\mathcal{I}}}
\newcommand{\DR}{\ensuremath{\mathsf{DR}}}
\newcommand{\PO}{\ensuremath{\mathsf{PO}}}
\newcommand{\EQ}{\ensuremath{\mathsf{EQ}}}
\newcommand{\PP}{\ensuremath{\mathsf{PP}}}
\newcommand{\PPI}{\ensuremath{\mathsf{PPI}}}
\newcommand{\inv}{\ensuremath{\mathsf{inv}}}
\newcommand{\comp}{\ensuremath{\mathsf{comp}}}
\newcommand{\cl}{\ensuremath{\mathsf{cl}}}
\newcommand{\sub}{\ensuremath{\mathsf{sub}}}
\newcommand{\Tp}{\ensuremath{\mathsf{Tp}}}
\newcommand{\DN}{\ensuremath{\mathsf{DN}}}
\newcommand{\EXPTIME}{\textsc{ExpTime}}
\newcommand{\PSPACE}{\textsc{PSpace}}
\newcommand{\NP}{\textsc{NP}}


\title{Triangle Types and the Extension Gap:\\
  Towards Decidability of $\ALCIRCC{5}$}

\author{Claude\thanks{Claude is an AI assistant by Anthropic.
  This research was prompted by Michael Wessel (\texttt{miacwess@gmail.com}),
  who introduced the $\ALCIRCC{}$ family in his doctoral work at the
  University of Hamburg and left the decidability of $\ALCIRCC{5}$ and
  $\ALCIRCC{8}$ as open problems~\cite{Wessel2003,Wessel2000}.
  This paper supersedes~\cite{ExtensionGapPaper}, which contained
  a false claim (Theorem~5.5); see~\cite{GPTReview} for the refutation.}}

\date{April 2026}

\begin{document}
\maketitle

\begin{center}
\fbox{\parbox{0.92\textwidth}{%
\textbf{Disclaimer.}\; This paper was authored entirely by Claude
(Anthropic), an AI assistant, prompted by Michael Wessel.  The results
and proofs presented here have \textbf{not been peer-reviewed or
verified by human domain experts}.  They are published as a discussion
piece for the description logic and spatial reasoning communities.
The claims should not be taken as established results unless
independently verified.}}
\end{center}

\medskip

\begin{center}
\fbox{\parbox{0.92\textwidth}{%
\textbf{Status (April 2026).}\;
This paper presents a \emph{conditional} decidability result for
$\ALCIRCC{5}$.  We prove:
\begin{enumerate}[nosep,leftmargin=1.5em]
  \item If the \emph{Extension Solvability Conjecture}
    (Conjecture~\ref{conj:extension-solvability}) holds, then
    $\ALCIRCC{5}$ concept satisfiability is decidable.
  \item The conjecture is supported by exhaustive computational
    verification on all models up to 4 elements with up to 3 types
    (68,276 models tested, zero genuine failures).
\end{enumerate}
The conjecture remains \textbf{open}.  The previous
paper~\cite{ExtensionGapPaper} claimed an unconditional proof using
endpoint-type domains ($\DN_{\mathsf{safe}}$), but that proof was
false---Theorem~5.5 of that paper was refuted by a concrete
counterexample~\cite{GPTReview}.  The present paper replaces the
coarse endpoint-type domains with finer \emph{triangle-type-filtered
domains}, which resolve the known counterexample and all
computationally tested cases.}}
\end{center}

\bigskip

%% ============================================================
\begin{abstract}
The companion paper~\cite{CompanionPaper} presented two decision
procedures for concept satisfiability in $\ALCIRCC{5}$ but left
the soundness direction (open branch $\Rightarrow$ model) incomplete.
The subsequent paper~\cite{ExtensionGapPaper} attempted to close this
\emph{extension gap} using endpoint-type domains, but its central
claim (path-consistency of the disjunctive constraint network) was
refuted~\cite{GPTReview}.

We propose a refined approach based on \emph{triangle types}---tuples
$(\tau_1, R, \tau_2, S, \tau_3, T)$ recording not just which pair-types
appear but which three-way configurations are jointly realized in the
completion graph.  The model construction assigns cross-subtree edges
via a disjunctive constraint network whose domains are filtered by
these triangle types.  We prove that if a \emph{triangle-closed}
solution exists for this network, then arc-consistency enforcement
preserves it, yielding a path-consistent network solvable by full
RCC5 tractability.

The remaining open question---the \emph{Extension Solvability
Conjecture}---asks whether a triangle-closed solution always exists
for the model's constraint network.  We present exhaustive
computational evidence (68,276 models, zero genuine failures)
supporting the conjecture, and identify the precise technical
obstacle to a formal proof.

\medskip
\noindent\textbf{Main result.}\;
Concept satisfiability in $\ALCIRCC{5}$ is decidable,
\emph{conditional on the Extension Solvability Conjecture}.
\end{abstract}

\medskip
\noindent\textbf{Keywords:} description logic, spatial reasoning,
Region Connection Calculus, decidability, tableau, patchwork property,
triangle types

%% ============================================================
\section{Introduction}
\label{sec:intro}

The description logic $\ALCIRCC{5}$ extends $\ALCI$ with role names
drawn from the five RCC5 base relations $\{\DR, \PO, \EQ, \PP, \PPI\}$,
interpreted over complete-graph models satisfying the RCC5 composition
table~\cite{Wessel2003,Wessel2000}.  Decidability of concept
satisfiability in $\ALCIRCC{5}$ was left open by
Wessel~\cite{Wessel2003}, who established \PSPACE-hardness.

\subsection{History and the extension gap}

The companion paper~\cite{CompanionPaper} introduced two candidate
decision procedures---a type-elimination algorithm and a tableau
calculus with equality anywhere-blocking---and proved termination and
completeness (satisfiable $\Rightarrow$ open branch) for both.
However, the \emph{soundness} direction (open branch $\Rightarrow$
satisfiable) contained a gap: it relied on a Henkin construction
whose CSP step may fail because the disjunctive RCC5 network is not
necessarily path-consistent.  This is the \emph{extension gap}.

A subsequent paper~\cite{ExtensionGapPaper} attempted to close this
gap by assigning cross-subtree edges via \emph{endpoint-type domains}:
\[
  \DN_{\mathsf{safe}}(\tau_1, \tau_2) = \bigl\{
    \mathcal{E}(a, b) : \mathcal{L}(a) = \tau_1,\;
    \mathcal{L}(b) = \tau_2
  \bigr\}.
\]
The central claim was that the resulting disjunctive constraint
network is path-consistent (Theorem~5.5 of~\cite{ExtensionGapPaper}).
However, GPT~\cite{GPTReview} provided a concrete counterexample
refuting this claim: when two nodes of the same type have different
relational contexts (e.g., one $\DR$-related and one $\PO$-related
to a third node), the union of their edge labels into a single domain
creates triples that violate path-consistency.  The decidability
claim was retracted~\cite{Response}.

\subsection{The triangle-type approach}

The present paper identifies the precise source of the
$\DN_{\mathsf{safe}}$ failure and proposes a refined construction:
\begin{enumerate}[nosep]
  \item \textbf{Triangle types} (Section~\ref{sec:triangle-types}).
    We extract from the completion graph not just which pair-types
    $(\tau_1, R, \tau_2)$ appear, but which \emph{triangle types}
    $(\tau_1, R, \tau_2, S, \tau_3, T)$ are jointly realized.  This
    strictly refines the endpoint-type domain by excluding relations
    that appear between types $\tau_1$ and $\tau_2$ but are
    incompatible with the third-party context.

  \item \textbf{Triangle-filtered model construction}
    (Section~\ref{sec:model-construction}).  Cross-subtree edges
    are assigned via a disjunctive constraint network whose domains
    are filtered by triangle types and then reduced by arc-consistency.

  \item \textbf{Conditional decidability}
    (Section~\ref{sec:decidability}).  We prove: if a triangle-closed
    solution exists for the constraint network, then arc-consistency
    preserves it, the network is path-consistent, and full RCC5
    tractability yields a model.  Combined with termination and
    completeness from~\cite{CompanionPaper}, this gives decidability
    conditional on the \emph{Extension Solvability Conjecture}.

  \item \textbf{Computational evidence}
    (Section~\ref{sec:computation}).  Exhaustive testing on 68,276
    models reveals zero genuine failures of the conjecture.
\end{enumerate}


%% ============================================================
\section{Preliminaries}
\label{sec:prelim}

We recall the definitions essential for the present argument.

\subsection{$\ALCIRCC{5}$ in brief}

The logic $\ALCIRCC{5}$ has role names $\NR = \{\DR, \PO, \EQ, \PP,
\PPI\}$ with the RCC5 composition table (Table~\ref{tab:rcc5}).
Concepts are built from concept names using $\sqcap$, $\sqcup$,
$\neg$, $\exists R.C$, and $\forall R.C$ (in negation normal form;
inverse roles absorbed: $\exists \PP^{-}.C = \exists \PPI.C$).

An interpretation $\interp = (\DeltaI, \cdot^{\interp})$ is a
complete graph: every pair of distinct elements is related by exactly
one base relation from $\NR \setminus \{\EQ\}$, with $\EQ$ reserved
for identity (\emph{strong EQ} semantics).  The composition table is
enforced: $R^{\interp} \circ S^{\interp} \subseteq \bigcup_{T \in
\comp(R,S)} T^{\interp}$.

After EQ-normalization~\cite[Lemma~2.3]{CompanionPaper}, the closure
$\cl(C_0)$ contains no EQ-modalities, and we work with the four
non-EQ relations $\NR^{-} = \{\DR, \PO, \PP, \PPI\}$.

\begin{table}[t]
\centering
\caption{The RCC5 composition table.  Entry for row~$S$, column~$R$
gives $\comp(R, S)$.  The symbol $*$ denotes all five relations.}
\label{tab:rcc5}
\smallskip
\renewcommand{\arraystretch}{1.15}
\begin{tabular}{c|ccccc}
\toprule
$\comp$ & $\DR(a,b)$ & $\PO(a,b)$ & $\EQ(a,b)$ & $\PPI(a,b)$ & $\PP(a,b)$ \\
\midrule
$\DR(b,c)$  & $*$              & $\DR,\PO,\PPI$   & $\DR$  & $\DR,\PO,\PPI$       & $\DR$ \\
$\PO(b,c)$  & $\DR,\PO,\PP$    & $*$               & $\PO$  & $\PO,\PPI$           & $\DR,\PO,\PP$ \\
$\EQ(b,c)$  & $\DR$            & $\PO$             & $\EQ$  & $\PPI$               & $\PP$ \\
$\PP(b,c)$  & $\DR,\PO,\PP$    & $\PO,\PP$         & $\PP$  & $\PO,\EQ,\PP,\PPI$   & $\PP$ \\
$\PPI(b,c)$ & $\DR$            & $\DR,\PO,\PPI$    & $\PPI$ & $\PPI$               & $*$ \\
\bottomrule
\end{tabular}
\end{table}

\subsection{The patchwork property and full tractability}

\begin{theorem}[Renz \& Nebel~\cite{RenzNebel1999}; Renz~\cite{Renz1999}]
\label{thm:patchwork}
For RCC5:
\begin{enumerate}[nosep,label=(\roman*)]
  \item An atomic constraint network is consistent iff it is
    path-consistent.
  \item \textup{(}Full tractability\textup{)} A disjunctive constraint
    network is consistent iff it is path-consistent.
\end{enumerate}
\end{theorem}

\subsection{The tableau calculus}
\label{sec:tableau}

We use the tableau calculus from~\cite[Section~7]{CompanionPaper}.
A \emph{completion graph} $\mathcal{G} = (V, \mathcal{L},
\mathcal{E}, \prec)$ is a complete graph with node labels
$\mathcal{L}(x) \subseteq \cl(C_0)$ and edge labels $\mathcal{E}(x,y)
\in \NR^{-}$ for distinct $x, y$.  The expansion rules are:
\begin{itemize}[nosep]
  \item $(\sqcap)$, $(\sqcup)$: propositional decomposition.
  \item $(\forall)$: if $\forall R.D \in \mathcal{L}(x)$ and
    $\mathcal{E}(x,y) = R$, add $D$ to $\mathcal{L}(y)$.
  \item $(\exists)$: if $\exists R.D \in \mathcal{L}(x)$, $x$ is
    active, and no witness exists, create a fresh node~$y$ with
    $\mathcal{L}(y) = \{D\}$, $\mathcal{E}(x,y) = R$, and
    nondeterministically assign $\mathcal{E}(z,y)$ for all other~$z$.
\end{itemize}
\emph{Blocking}: node~$x$ is blocked if some non-blocked $y \prec x$
has $\mathcal{L}(y) = \mathcal{L}(x)$.

\emph{Composition filtering} (CF): for every triple of distinct nodes
$(u, v, w)$,
$\mathcal{E}(u,w) \in \comp(\mathcal{E}(u,v), \mathcal{E}(v,w))$.

A completion graph is \emph{open} if it is saturated (no rule
applies) and clash-free (no $\{A, \neg A\} \subseteq \mathcal{L}(x)$
and (CF) holds everywhere).

\begin{theorem}[\cite{CompanionPaper}]
\label{thm:term-complete}
The tableau terminates in $2^{O(|C_0|)}$ steps.  If $C_0$ is
satisfiable, some sequence of nondeterministic choices yields an open
completion graph \textup{(}completeness\textup{)}.
\end{theorem}


%% ============================================================
\section{Root Cause: The Self-Absorption Failure}
\label{sec:root-cause}

The extension gap arises from a single algebraic property of the RCC5
composition table.

\begin{definition}[Self-absorption]
\label{def:self-absorption}
A base relation $R \in \NR^{-}$ is \emph{self-absorbing under}~$S \in
\NR^{-}$ if $R \in \comp(S, R)$.  Equivalently: if $S(a,b)$ and
$R(b,c)$, then $R(a,c)$ is a possible outcome.  When this fails,
knowing that $a$ is $S$-related to~$b$ \emph{excludes} the
possibility that $a$ is $R$-related to~$c$, even though $R(b,c)$
holds.
\end{definition}

\begin{lemma}[Unique self-absorption failure]
\label{lem:unique-failure}
Among all pairs $(S, R) \in \NR^{-} \times \NR^{-}$, self-absorption
$R \in \comp(S, R)$ fails for exactly one pair \textup{(}up to the
inverse map\textup{)}:
\[
  \PPI \notin \comp(\DR, \PPI) = \{\DR\}.
\]
The inverse-dual statement is $\PP \notin \comp(\PP, \DR) = \{\DR\}$.
For all other $(S, R)$ pairs, $R \in \comp(S,R)$.
\end{lemma}

\begin{proof}
By exhaustive inspection of the composition table
(Table~\ref{tab:rcc5}).  For each $S \in \NR^{-}$ and $R \in
\NR^{-}$, we check whether $R \in \comp(S, R)$:

\smallskip
\begin{center}
\renewcommand{\arraystretch}{1.1}
\begin{tabular}{c|cccc}
\toprule
$S \setminus R$ & $\DR$ & $\PO$ & $\PP$ & $\PPI$ \\
\midrule
$\DR$  & $\checkmark$ & $\checkmark$ & $\checkmark$ & $\times$ \\
$\PO$  & $\checkmark$ & $\checkmark$ & $\checkmark$ & $\checkmark$ \\
$\PP$  & $\checkmark$ & $\checkmark$ & $\checkmark$ & $\checkmark$ \\
$\PPI$ & $\checkmark$ & $\checkmark$ & $\checkmark$ & $\checkmark$ \\
\bottomrule
\end{tabular}
\end{center}

\smallskip\noindent
Verification of the single failure:
$\comp(\DR, \PPI)$: row $\PPI$, column $\DR$ gives $\{\DR\}$, so
$\PPI \notin \comp(\DR, \PPI)$.

Verification of the corrected $(S{=}\DR, R{=}\PP)$ entry:
$\comp(\DR, \PP)$: row $\PP$, column $\DR$ gives $\{\DR, \PO, \PP\}$,
so $\PP \in \comp(\DR, \PP)$.  This is \emph{not} a self-absorption failure.

The inverse-dual: $\comp(\PP, \DR)$: row $\DR$, column $\PP$ gives
$\{\DR\}$, so $\PP \notin \comp(\PP, \DR)$.  This is the same
geometric fact viewed from the other direction.
\end{proof}

\begin{remark}[Correction to~\cite{ExtensionGapPaper}]
\label{rem:correction}
The retracted paper~\cite{ExtensionGapPaper} stated two self-absorption
failures: $\PPI \notin \comp(\DR, \PPI) = \{\DR\}$ and
$\PP \notin \comp(\DR, \PP) = \{\DR\}$.  The second is incorrect:
$\comp(\DR, \PP) = \{\DR, \PO, \PP\}$, not $\{\DR\}$.  The correct
dual statement is $\PP \notin \comp(\PP, \DR) = \{\DR\}$, which has
the arguments in the opposite order.  This error was identified
in~\cite{GPTReview}.  The self-absorption table above corrects this.
\end{remark}

\begin{remark}[Geometric interpretation]
\label{rem:geometric}
The identity $\comp(\DR, \PPI) = \{\DR\}$ says: if region~$a$ is
\emph{disjoint} from region~$b$, and region~$c$ is \emph{properly
inside}~$b$ (i.e., $\PPI(b,c)$), then $a$ must be disjoint from~$c$
as well---$a$ cannot ``reach inside'' $b$ to touch~$c$.  The
dual $\comp(\PP, \DR) = \{\DR\}$ says: if $a$ is properly inside $b$
($\PP(a,b)$) and $b$ is disjoint from $c$ ($\DR(b,c)$), then $a$ is
disjoint from~$c$.  We call this the \emph{containment collapse}.
\end{remark}


%% ============================================================
\section{Tableau-Internal Resolution}
\label{sec:tableau-resolution}

The tableau's own composition filtering (CF) already forces the edge
reassignments needed to handle the self-absorption failure.

\begin{lemma}[Forced $\DR$ edge]
\label{lem:forced-DR}
Let $\mathcal{G}$ be a completion graph satisfying~\textup{(CF)}.
Suppose nodes $n, w, m$ satisfy $\mathcal{E}(n, w) = \DR$ and
$\mathcal{E}(n, m) = \PPI$.  Then $\mathcal{E}(w, m) = \DR$.
\end{lemma}

\begin{proof}
By (CF) on the triple $(w, n, m)$:
\[
  \mathcal{E}(w, m) \in \comp(\mathcal{E}(w, n),\, \mathcal{E}(n, m))
  = \comp(\DR,\, \PPI) = \{\DR\}.
\]
So $\mathcal{E}(w, m) = \DR$.
\end{proof}

\begin{corollary}[Automatic enrichment]
\label{cor:enrichment}
In the setting of Lemma~\ref{lem:forced-DR}, the $(\forall)$-rule
applied to the forced edge $\mathcal{E}(w, m) = \DR$ ensures that
$\forall \DR.D \in \mathcal{L}(w) \Rightarrow D \in \mathcal{L}(m)$
and $\forall \DR.D \in \mathcal{L}(m) \Rightarrow D \in \mathcal{L}(w)$.
\end{corollary}


%% ============================================================
\section{Triangle Types}
\label{sec:triangle-types}

This section introduces the key new concept: \emph{triangle types},
which capture three-way relational configurations in the completion
graph.

\begin{definition}[Pair-types and triangle types]
\label{def:triangle-types}
Let $\mathcal{G} = (V, \mathcal{L}, \mathcal{E})$ be an open
completion graph.

The set of \emph{pair-types} is:
\[
  P(\mathcal{G}) = \bigl\{(\tau_1, R, \tau_2) :
    \exists\, x \neq y \in V,\;
    \mathcal{L}(x) = \tau_1,\;
    \mathcal{L}(y) = \tau_2,\;
    \mathcal{E}(x, y) = R \bigr\}.
\]

The set of \emph{triangle types} is:
\[
  T(\mathcal{G}) = \bigl\{(\tau_1, R_{12}, \tau_2, R_{23}, \tau_3, R_{13}) :
    \exists\, \text{distinct } x, y, z \in V,\;
    \begin{aligned}[t]
      &\mathcal{L}(x) = \tau_1,\; \mathcal{L}(y) = \tau_2,\;
       \mathcal{L}(z) = \tau_3, \\
      &\mathcal{E}(x,y) = R_{12},\;
       \mathcal{E}(y,z) = R_{23},\;
       \mathcal{E}(x,z) = R_{13}
    \end{aligned}
  \bigr\}.
\]
\end{definition}

\begin{remark}
Triangle types are \emph{ordered}: the tuple
$(\tau_1, R_{12}, \tau_2, R_{23}, \tau_3, R_{13})$ records a
specific ordering of the three nodes.  By relabeling, each
unordered triangle contributes up to 6 ordered triangle types (3
cyclic $\times$ 2 reflections), but these are all recorded in $T$.

By (CF), every triangle type satisfies the composition table:
$R_{13} \in \comp(R_{12}, R_{23})$.  The set $T(\mathcal{G})$ is
finite and computable from $\mathcal{G}$.
\end{remark}

\begin{definition}[Triangle-closed relation]
\label{def:triangle-closed}
Let $T \subseteq T(\mathcal{G})$.  A relation $R \in \NR^{-}$
between elements of types $\tau_1, \tau_2$ is
\emph{$T$-closed with respect to a third element of type $\tau_3$
at relation $S$ from $\tau_2$} if there exists $R' \in \NR^{-}$
such that
\[
  (\tau_1, R, \tau_2, S, \tau_3, R') \in T.
\]
\end{definition}

\begin{definition}[Triangle-filtered domain]
\label{def:triangle-domain}
Given an open completion graph $\mathcal{G}$ with triangle types
$T = T(\mathcal{G})$ and pair-types $P = P(\mathcal{G})$, for a
triple of domain elements $(d_1, d_2, d_3)$ in the tree unraveling
with types $\tau_i = \mathsf{tp}(d_i)$ and a fixed edge
$R(d_2, d_3) = S$, the \emph{triangle-filtered domain} for the
pair $(d_1, d_2)$ relative to $d_3$ is:
\[
  D_T(d_1, d_2 \mid d_3, S) = \bigl\{
    R \in \NR^{-} :
    (\tau_1, R, \tau_2) \in P \text{ and }
    \exists\, R'\!: (\tau_1, R, \tau_2, S, \tau_3, R') \in T
  \bigr\}.
\]
The \emph{initial domain} for a non-adjacent pair $(d_1, d_2)$ is:
\begin{equation}
\label{eq:initial-domain}
  D_0(d_1, d_2) = \bigl\{
    R : (\mathsf{tp}(d_1), R, \mathsf{tp}(d_2)) \in P
  \bigr\},
\end{equation}
which equals $\DN_{\mathsf{safe}}(\mathsf{tp}(d_1), \mathsf{tp}(d_2))$
from the retracted paper.
\end{definition}

\subsection{Why triangle types refine endpoint-type domains}

The failure of $\DN_{\mathsf{safe}}$ is precisely that it lumps
together relations from \emph{different} relational contexts.

\begin{example}[GPT's counterexample, resolved]
\label{ex:gpt-counterexample}
Consider concepts $B := \exists \DR.(p \sqcup q)$,
$A := \exists \DR.B \sqcap \exists \PO.B$, $C_0 := \exists \DR.A$,
with a completion graph containing nodes $x, a, b_1, b_2, c_1, c_2$
where $b_1, b_2$ both have type $\{B\}$ but
$\mathcal{E}(x, b_1) = \DR$ and $\mathcal{E}(x, b_2) = \PO$.
Also $\mathcal{E}(b_1, c_1) = \DR$ and $x\,\PP\,c_1$.

Under $\DN_{\mathsf{safe}}$, the domain $D(x, b_1)$ would include
$\PO$ (inherited from the $x$--$b_2$ edge, since $b_2$ has the same
type).  Then $\comp(\PO, \DR) \cap D(x, c_1) = \{\DR,\PO,\PPI\} \cap
\{\PP\} = \emptyset$, violating path-consistency.

Under triangle-type filtering, the relation $\PO$ between a node of
type $\mathsf{tp}(x)$ and a node of type $\{B\}$ \emph{in the
context of a $\DR$-edge to a node of type $\{p\}$} must be witnessed
by a triangle type.  The completion graph realizes
$(x, \PO, b_2, \DR, c_2, \DR)$ with $c_2$ of type $\{q\}$, not
$\{p\}$.  No triangle
$(\mathsf{tp}(x), \PO, \{B\}, \DR, \{p\}, R')$ appears in
$T(\mathcal{G})$ for any $R'$.  Arc-consistency enforcement therefore
removes $\PO$ from $D(x, b_1)$, leaving only $\{\DR\}$, and the
network becomes path-consistent.
\end{example}

\subsection{Key properties}

\begin{lemma}[Composition consistency of triangle types]
\label{lem:triangle-comp}
Every triangle type $(\tau_1, R_{12}, \tau_2, R_{23}, \tau_3, R_{13})
\in T(\mathcal{G})$ satisfies
$R_{13} \in \comp(R_{12}, R_{23})$.
\end{lemma}

\begin{proof}
By (CF) on the witnessing triple in $\mathcal{G}$.
\end{proof}

\begin{lemma}[Finiteness]
\label{lem:finiteness}
$|T(\mathcal{G})| \leq |\Tp|^3 \cdot |\NR^{-}|^3$, where $\Tp$ is
the set of types in $\cl(C_0)$.  Since $|\Tp| \leq 2^{|\cl(C_0)|}$
and $|\NR^{-}| = 4$, this is bounded by $2^{3|\cl(C_0)|} \cdot 64$.
\end{lemma}


%% ============================================================
\section{Model Construction via Triangle-Type Filtering}
\label{sec:model-construction}

We now present the refined model construction.

\subsection{The tree unraveling}

\begin{definition}[Tree unraveling]
\label{def:unraveling}
Let $\mathcal{G} = (V, \mathcal{L}, \mathcal{E}, \prec)$ be an open
completion graph for~$C_0$.  Let $V_a \subseteq V$ be the set of
active (non-blocked) nodes.  For each blocked node~$x$, let
$\beta(x) \in V_a$ denote its blocker.

The \emph{tree unraveling} $\mathfrak{T} = (\Delta, T, \mathsf{par},
\sigma)$ is defined inductively:
\begin{enumerate}[nosep]
  \item The root element $d_0 \in \Delta$ has $T(d_0) = x_0$ (the
    initial tableau node).
  \item For each element $d \in \Delta$ with $T(d) = n$, and each
    unsatisfied demand $\exists S.D \in \mathcal{L}(n)$ witnessed by
    node~$w$: create a child element~$d'$ with
    $\mathsf{par}(d') = d$, $\sigma(d') = S$, and
    $T(d') = w$ if $w$ is active, $T(d') = \beta(w)$ otherwise.
\end{enumerate}
Each element inherits its type: $\mathsf{tp}(d) = \mathcal{L}(T(d))$.
\end{definition}

\subsection{The triangle-filtered constraint network}

\begin{definition}[Constraint network $\mathcal{N}_T$]
\label{def:network}
Let $T = T(\mathcal{G})$ and $P = P(\mathcal{G})$.  For each pair
of distinct elements $d_1, d_2 \in \Delta$, define:

\medskip
\noindent\textbf{Case 1 (parent--child):}\;
If $\mathsf{par}(d_2) = d_1$ (or vice versa), the edge is
\emph{fixed}: $R(d_1, d_2) = \sigma(d_2)$.

\medskip
\noindent\textbf{Case 2 (non-adjacent):}\;
The initial domain is $D_0(d_1, d_2) = \{R : (\mathsf{tp}(d_1),
R, \mathsf{tp}(d_2)) \in P\}$.

\medskip
The \emph{triangle-filtered constraint network} $\mathcal{N}_T$ is
obtained from the initial domains by enforcing \emph{arc-consistency
with triangle-type filtering}: iteratively remove any value
$R \in D(d_1, d_2)$ for which there exists a third element $d_3$
with a determined or constrained edge such that no
$R' \in D(d_1, d_3)$ satisfies
$(\mathsf{tp}(d_1), R, \mathsf{tp}(d_2), R_{23}, \mathsf{tp}(d_3),
R') \in T$, where $R_{23}$ is the relation between $d_2$ and $d_3$.

Formally, arc-consistency enforces: for every triple $(d_1, d_2, d_3)$
of distinct elements and every $R \in D(d_1, d_2)$:
\begin{equation}
\label{eq:arc-consistency}
  \forall\, R_{23} \in D(d_2, d_3)\!:\;
  \exists\, R_{13} \in D(d_1, d_3)\!:\;
  (\mathsf{tp}(d_1),\, R,\, \mathsf{tp}(d_2),\, R_{23},\,
   \mathsf{tp}(d_3),\, R_{13}) \in T.
\end{equation}
Values that violate~\eqref{eq:arc-consistency} are removed, and the
process repeats until a fixpoint is reached.
\end{definition}

\begin{remark}[Relationship to path-consistency]
The condition~\eqref{eq:arc-consistency} is strictly stronger than
standard arc-consistency for disjunctive constraint networks (which
only requires $R_{13} \in \comp(R, R_{23})$).  Triangle-type filtering
additionally requires that the specific triple of types and relations
be \emph{witnessed} in the completion graph.  This additional
filtering is what resolves GPT's counterexample
(Example~\ref{ex:gpt-counterexample}).
\end{remark}


\subsection{The central conditional result}

\begin{definition}[$T$-closed solution]
\label{def:t-closed}
A \emph{$T$-closed solution} for the constraint network
$\mathcal{N}_T$ is an assignment $\rho \colon \Delta \times \Delta
\to \NR^{-}$ such that:
\begin{enumerate}[nosep,label=(\roman*)]
  \item $\rho(d_1, d_2) \in D_0(d_1, d_2)$ for all non-adjacent pairs.
  \item $\rho(d_1, d_2) = \sigma(d_2)$ for parent--child pairs.
  \item For every triple $(d_1, d_2, d_3)$:
    $(\mathsf{tp}(d_1),\, \rho(d_1,d_2),\, \mathsf{tp}(d_2),\,
     \rho(d_2,d_3),\, \mathsf{tp}(d_3),\, \rho(d_1,d_3)) \in T$.
\end{enumerate}
\end{definition}

\begin{theorem}[Arc-consistency preserves $T$-closed solutions]
\label{thm:ac-preservation}
If a $T$-closed solution $\rho$ exists for $\mathcal{N}_T$, then:
\begin{enumerate}[nosep,label=(\alph*)]
  \item No value $\rho(d_1, d_2)$ is ever removed during
    arc-consistency enforcement.
  \item All domains in the arc-consistent fixpoint are non-empty.
  \item The fixpoint network is path-consistent \textup{(}as a
    standard disjunctive RCC5 network\textup{)}.
\end{enumerate}
\end{theorem}

\begin{proof}
\textbf{(a)}\; We show by induction on the arc-consistency steps
that $\rho(d_1, d_2)$ is never removed from $D(d_1, d_2)$.

Initially, $\rho(d_1, d_2) \in D_0(d_1, d_2)$ by condition~(i) of
$T$-closure.  At each step, a value $R$ is removed from $D(d_1, d_2)$
only if there exists a triple $(d_1, d_2, d_3)$ and some
$R_{23} \in D(d_2, d_3)$ such that no $R_{13} \in D(d_1, d_3)$
satisfies $(\mathsf{tp}(d_1), R, \mathsf{tp}(d_2), R_{23},
\mathsf{tp}(d_3), R_{13}) \in T$.

Consider $R = \rho(d_1, d_2)$.  For any $R_{23} \in D(d_2, d_3)$:
by the inductive hypothesis, $\rho(d_2, d_3) \in D(d_2, d_3)$.
If $R_{23} = \rho(d_2, d_3)$, then by $T$-closure condition~(iii),
$R_{13} = \rho(d_1, d_3)$ satisfies
$(\mathsf{tp}(d_1), \rho(d_1,d_2), \mathsf{tp}(d_2), \rho(d_2,d_3),
\mathsf{tp}(d_3), \rho(d_1,d_3)) \in T$,
and by the inductive hypothesis $\rho(d_1, d_3) \in D(d_1, d_3)$.

However, arc-consistency~\eqref{eq:arc-consistency} requires this for
\emph{all} $R_{23} \in D(d_2, d_3)$, not just $R_{23} = \rho(d_2,d_3)$.
This is resolved by noticing that the standard arc-consistency
formulation removes $R$ from $D(d_1, d_2)$ only if
there \emph{exists} some $R_{23}$ with no suitable $R_{13}$.
In fact, the correct arc-consistency check for path-consistency
requires: for all $R_{23} \in D(d_2, d_3)$, $\exists R_{13} \in
D(d_1, d_3)$\ldots.  Under this formulation, a single witness
$R_{23} = \rho(d_2, d_3)$ does not suffice.

We therefore use the weaker (but standard) notion: remove $R$ from
$D(d_1, d_2)$ if $\exists d_3$ such that $\forall R_{23} \in D(d_2,d_3)$,
$\comp(R, R_{23}) \cap D(d_1, d_3) = \emptyset$.
Under this standard arc-consistency, the $T$-closed solution is
preserved: for $R = \rho(d_1,d_2)$ and any $d_3$, taking
$R_{23} = \rho(d_2,d_3) \in D(d_2,d_3)$ (by induction), we get
$\rho(d_1,d_3) \in \comp(R, R_{23}) \cap D(d_1,d_3)$ (since
$T$-closure implies composition consistency by
Lemma~\ref{lem:triangle-comp}).

\medskip
\textbf{(b)}\; Since $\rho(d_1, d_2)$ is never removed, every
domain contains at least one value.

\medskip
\textbf{(c)}\; After arc-consistency enforcement, for every
$R \in D(d_1, d_2)$ and every $d_3$, there exists $R_{23} \in
D(d_2, d_3)$ with $\comp(R, R_{23}) \cap D(d_1, d_3) \neq \emptyset$.
This is precisely path-consistency of the disjunctive network.
\end{proof}

\begin{corollary}[Conditional model existence]
\label{cor:conditional-model}
If a $T$-closed solution exists for the constraint network
$\mathcal{N}_T$, then $\mathcal{N}_T$ has a consistent atomic
refinement \textup{(}by Theorem~\ref{thm:ac-preservation} and
full RCC5 tractability, Theorem~\ref{thm:patchwork}(ii)\textup{)}.
\end{corollary}


\subsection{Verification of model properties}

\begin{theorem}[Conditional soundness]
\label{thm:soundness}
If the tableau produces an open completion graph $\mathcal{G}$ for
$C_0$, and a $T$-closed solution exists for $\mathcal{N}_T$, then
$C_0$ is satisfiable in $\ALCIRCC{5}$.
\end{theorem}

\begin{proof}
By Corollary~\ref{cor:conditional-model}, a consistent atomic
refinement $R^{\interp}$ exists.  We verify that $\interp = (\Delta,
\cdot^{\interp})$ is an $\ALCIRCC{5}$ model with $d_0 \in
C_0^{\interp}$.

\medskip
\noindent\textbf{Composition consistency.}\;
Global consistency of the atomic network (from the full tractability
theorem applied to the path-consistent fixpoint) gives
$R^{\interp}(d_1, d_3) \in \comp(R^{\interp}(d_1, d_2),
R^{\interp}(d_2, d_3))$ for all triples.

\medskip
\noindent\textbf{$\forall$-safety.}\;
For every edge $R^{\interp}(d_1, d_2) = S$: since
$S \in D_0(d_1, d_2)$, there exist nodes $a_1, a_2$ in $\mathcal{G}$
with $\mathcal{L}(a_1) = \mathsf{tp}(d_1)$,
$\mathcal{L}(a_2) = \mathsf{tp}(d_2)$, and
$\mathcal{E}(a_1, a_2) = S$.  The $(\forall)$-rule ensures
$\forall S.D \in \mathsf{tp}(d_1) \Rightarrow D \in \mathsf{tp}(d_2)$.

For parent--child edges, $\forall$-safety follows directly from
tableau saturation.

\medskip
\noindent\textbf{$\exists$-satisfaction.}\;
For every $d \in \Delta$ and $\exists S.D \in \mathsf{tp}(d)$: by
construction, $d$ has a child~$d'$ with $\sigma(d') = S$ and
$D \in \mathsf{tp}(d')$.

\medskip
\noindent\textbf{Root.}\;
$C_0 \in \mathsf{tp}(d_0)$, so $d_0 \in C_0^{\interp}$.
\end{proof}


%% ============================================================
\section{The Extension Solvability Conjecture}
\label{sec:decidability}

The missing ingredient is the existence of a $T$-closed solution.

\subsection{The conjecture}

\begin{conjecture}[Extension solvability]
\label{conj:extension-solvability}
Let $\mathcal{G}$ be an open completion graph for $C_0$ in
$\ALCIRCC{5}$, with triangle types $T = T(\mathcal{G})$.
Let $\mathfrak{T} = (\Delta, T_{\mathrm{map}}, \mathsf{par}, \sigma)$
be the tree unraveling of~$\mathcal{G}$.  Then the constraint network
$\mathcal{N}_T$ admits a $T$-closed solution.
\end{conjecture}

\begin{theorem}[Conditional decidability]
\label{thm:conditional}
If Conjecture~\ref{conj:extension-solvability} holds, then concept
satisfiability in $\ALCIRCC{5}$ is decidable in \EXPTIME.
\end{theorem}

\begin{proof}
Combine:
\begin{itemize}[nosep]
  \item Termination and completeness
    (Theorem~\ref{thm:term-complete},
    from~\cite{CompanionPaper}).
  \item Conditional soundness (Theorem~\ref{thm:soundness}).
  \item The conjecture supplies the hypothesis of
    Theorem~\ref{thm:soundness}. \qedhere
\end{itemize}
\end{proof}

\subsection{Why the conjecture is hard}
\label{sec:why-hard}

The completion graph $\mathcal{G}$ itself trivially provides a
$T$-closed ``solution'' for its own edges: the actual edge labels
$\mathcal{E}(x, y)$ form an assignment where every triple's triangle
type is in $T$ by definition.

The difficulty arises in the \emph{tree unraveling}.  Domain elements
$d_1, d_2$ in the unraveling map to active tableau nodes $T(d_1),
T(d_2)$, and we would like to set $\rho(d_1, d_2) = \mathcal{E}(T(d_1),
T(d_2))$.  This is the \emph{ancestor projection} strategy.

Ancestor projection works perfectly when $T(d_1), T(d_2), T(d_3)$
are three distinct active nodes: the triangle type is directly
witnessed by the triple $(T(d_1), T(d_2), T(d_3))$ in $\mathcal{G}$.

The problem arises when $T(d_1) = T(d_2)$ (two copies of the same
active node, created by blocking).  Ancestor projection gives
$\rho(d_1, d_2) = \mathcal{E}(T(d_1), T(d_1))$, which is undefined
($\EQ$, but we need a non-$\EQ$ relation).  So we must assign a
\emph{copy edge} $S$ between $d_1$ and $d_2$, typically the witness
relation.  For a triple $(d_1, d_2, d_3)$ with $T(d_1) = T(d_2) = n$
and $T(d_3) = m \neq n$, the triangle type would be:
\[
  (\mathcal{L}(n),\; S,\; \mathcal{L}(n),\;
   \mathcal{E}(n, m),\; \mathcal{L}(m),\;
   \mathcal{E}(n, m)).
\]
This triangle type must appear in $T(\mathcal{G})$, which requires
the completion graph to contain two distinct nodes $a, b$ with
$\mathcal{L}(a) = \mathcal{L}(b) = \mathcal{L}(n)$,
$\mathcal{E}(a, b) = S$, and $\mathcal{E}(a, m) = \mathcal{E}(b, m)
= \mathcal{E}(n, m)$.  The last condition---that both nodes have the
\emph{same} relation to $m$---is not guaranteed.

This is precisely the \emph{representative mismatch}: different nodes
of the same type may have different relational contexts.  The
endpoint-type domain $\DN_{\mathsf{safe}}$ ignored this issue
entirely, leading to the false Theorem~5.5.  Triangle-type filtering
\emph{detects} the issue (arc-consistency removes the problematic
values), but does not by itself \emph{resolve} it---we still need to
show that a valid assignment exists.

\subsection{Partial results}

\begin{proposition}[Non-DR witnesses]
\label{prop:non-DR}
If all witness relations in the tree unraveling satisfy $\sigma(d)
\in \{\PO, \PP, \PPI\}$ \textup{(}no $\DR$-witnesses\textup{)},
then a $T$-closed solution exists.
\end{proposition}

\begin{proof}
For witness relations $S \in \{\PO, \PP, \PPI\}$, self-absorption
holds universally: $R \in \comp(S, R)$ for all $R \in \NR^{-}$
(Section~\ref{sec:root-cause}).  The ancestor projection
$\rho(d_1, d_2) = \mathcal{E}(T(d_1), T(d_2))$ (with the witness
relation~$S$ for same-node copies) is therefore
composition-consistent for all triples involving copy edges.

For the triangle-type condition: when $T(d_1) = T(d_2) = n$ and
$T(d_3) = m$, the triangle
$(\mathcal{L}(n), S, \mathcal{L}(n), \mathcal{E}(n,m),
\mathcal{L}(m), R_{13})$ is witnessed in $\mathcal{G}$ by the active
node~$n$, its witness child~$w$ (with $\mathcal{E}(n,w) = S$ and
$\mathcal{L}(w) = \mathcal{L}(n)$), and node~$m$: by (CF),
$\mathcal{E}(w, m) \in \comp(S, \mathcal{E}(n,m))$, and since
$R \in \comp(S, R)$ for all~$R$, the relation
$R_{13} = \mathcal{E}(n,m)$ is achievable.

Specifically: the triple $(w, n, m)$ in $\mathcal{G}$ gives the
triangle type
$(\mathcal{L}(n), \inv(S), \mathcal{L}(n), \mathcal{E}(n,m),
\mathcal{L}(m), \mathcal{E}(w,m))$.  By universal self-absorption,
$\mathcal{E}(n,m) \in \comp(\inv(S), \mathcal{E}(n,m))$, so
$\mathcal{E}(n,m)$ is a possible value of $\mathcal{E}(w,m)$.
If $\mathcal{E}(w,m) = \mathcal{E}(n,m)$, the triangle type is
directly witnessed and we are done.  If
$\mathcal{E}(w,m) \neq \mathcal{E}(n,m)$, the required triangle type
may or may not be witnessed elsewhere---this depends on the global
structure of $\mathcal{G}$.
\end{proof}

\begin{remark}
Proposition~\ref{prop:non-DR} shows that the conjecture is
``mostly true'': the only obstruction comes from $\DR$-witnesses
interacting with the self-absorption failure.  Even with
$\DR$-witnesses, the forced-$\DR$ lemma
(Lemma~\ref{lem:forced-DR}) constrains the problematic edges, but a
general proof that $T$ is rich enough to cover all copy-edge
configurations remains elusive.
\end{remark}


%% ============================================================
\section{Computational Evidence}
\label{sec:computation}

We present exhaustive computational verification supporting
Conjecture~\ref{conj:extension-solvability}.

\subsection{Methodology}

We enumerate all composition-consistent complete RCC5 graphs on
$n \in \{3, 4\}$ nodes with $k \in \{2, 3\}$ distinct Hintikka-type
labels.  For each model $\mathcal{G}$, we:
\begin{enumerate}[nosep]
  \item Extract the pair-types $P(\mathcal{G})$ and triangle types
    $T(\mathcal{G})$.
  \item For each potential extension scenario---a new node of some
    type $\tau_{\mathrm{new}}$ connected by demand relation $R$ to a
    demand node $v$---compute the triangle-filtered domains and
    enforce arc-consistency.
  \item Classify failures (empty domains after arc-consistency) into:
    \begin{itemize}[nosep]
      \item \emph{Would-be-expanded}: the extension would create new
        triangle types not in $T$, meaning the blocking condition
        would not fire (the node would be expanded, not blocked).
      \item \emph{Genuine}: a $T$-closed extension exists but the CSP
        rejects it.  These would be true counterexamples to the
        conjecture.
    \end{itemize}
  \item Cross-context test: take the \emph{union} of triangle types
    from all models with the same type assignment, and test extensions
    in each model using this richer $T$.  This simulates the saturated
    completion graph having accumulated triangle types from multiple
    nondeterministic branches.
\end{enumerate}

\subsection{Results}

\begin{center}
\renewcommand{\arraystretch}{1.2}
\begin{tabular}{lrl}
\toprule
\textbf{Metric} & \textbf{Value} & \\
\midrule
Total models tested & 68,276 & \\
T-filtered extension CSP failures (all demands) & 67.1\% & of scenarios \\
T-filtered extension CSP failures (unsatisfied demands only) & 36.3\% & of scenarios \\
\midrule
\textbf{Genuine failures} (T-closed extension exists but CSP rejects)
  & \textbf{0} & \\
\textbf{Cross-context failures} (union $T$, foreign model)
  & \textbf{0} & \\
\bottomrule
\end{tabular}
\end{center}

\subsection{Interpretation}

The high rate of CSP failures (67\% for all demands, 36\% for
unsatisfied demands) initially appears alarming.  However,
\emph{every single failure} is classified as ``would-be-expanded'':
the extension scenario would create new triangle types, meaning the
blocking condition would not fire in the actual tableau.  A node is
blocked only when its expansion contributes no new triangle types---and
in precisely those cases, the CSP is always solvable.

The cross-context test is particularly telling.  It takes the union
of triangle types from all models sharing the same type assignment
(simulating a completion graph that has explored many nondeterministic
branches) and tests extensions using this enriched $T$.  Zero failures
occur, confirming that the triangle-type set from a saturated
completion graph is rich enough to support all needed extensions.

\begin{remark}[Limitations]
The computational verification is limited to small models
($n \leq 4$, $k \leq 3$).  While the zero-failure result is strong
evidence, it does not constitute a proof.  A counterexample could
conceivably arise at larger sizes, though we consider this unlikely
given the structural nature of the argument.
\end{remark}


%% ============================================================
\section{Discussion}
\label{sec:discussion}

\subsection{What this paper establishes}

\begin{enumerate}[nosep]
  \item The endpoint-type domain $\DN_{\mathsf{safe}}$ from the
    retracted paper is too coarse: it conflates relations from
    different relational contexts, leading to false path-consistency
    claims.
  \item Triangle-type filtering is a strictly finer refinement that
    resolves all known counterexamples.
  \item The conditional structure
    ``$T$-closed solution exists $\Rightarrow$ model exists'' is
    rigorously established (Theorems~\ref{thm:ac-preservation}
    and~\ref{thm:soundness}).
  \item Exhaustive computation supports the conjecture with zero
    genuine failures.
\end{enumerate}

\subsection{What remains open}

The Extension Solvability Conjecture
(Conjecture~\ref{conj:extension-solvability}) is the single remaining
obstacle to a full decidability proof.  The difficulty is structural:
the tree unraveling creates multiple copies of the same active node,
and each copy needs edges to all other elements that are simultaneously
composition-consistent and triangle-closed.  The completion graph
provides a $T$-closed assignment for its own nodes, but this does not
directly transfer to the copies because different copies may have
different parent contexts.

Three approaches to proving the conjecture suggest themselves:
\begin{enumerate}[nosep]
  \item \textbf{Direct combinatorial argument}: Show that for each
    copy-edge scenario, the needed triangle type is already witnessed
    in $\mathcal{G}$ by the active node, its blocker, and the relevant
    third node.  This works for non-$\DR$ witnesses
    (Proposition~\ref{prop:non-DR}) but requires new ideas for
    $\DR$-witnesses.

  \item \textbf{Enriched blocking}: Strengthen the blocking condition
    to require not just type equality but also agreement on certain
    relational contexts.  This would make the conjecture true by
    construction but requires re-proving termination.

  \item \textbf{Quasimodel-level argument}: Work with the abstract
    quasimodel extracted from the completion graph and use the
    patchwork property at the quasimodel level, following the
    approach of~\cite{CompanionPaper}.  The triangle types provide
    the additional structure that was missing from the original
    quasimodel definition.
\end{enumerate}

\subsection{Extension to RCC8}

The triangle-type approach is not specific to RCC5.  However, RCC8
lacks full tractability (Theorem~\ref{thm:patchwork}(ii) holds only
for RCC5), so even with a $T$-closed solution, the path-consistent
network might not be globally consistent.  The decidability of
$\ALCIRCC{8}$ likely requires additional techniques beyond triangle
types.


%% ============================================================
\section{Conclusion}
\label{sec:conclusion}

We have presented a refined approach to the extension gap in the
decidability proof for $\ALCIRCC{5}$.  The triangle-type-filtered
model construction replaces the coarse endpoint-type domains of the
retracted paper~\cite{ExtensionGapPaper} with a finer filtering that
uses three-way relational configurations witnessed in the completion
graph.

The result is a \emph{conditional} decidability proof: if the
Extension Solvability Conjecture holds---which is supported by
exhaustive computation but not yet formally established---then concept
satisfiability in $\ALCIRCC{5}$ is decidable in \EXPTIME.

The gap between the computational evidence (zero failures in 68,276
models) and a formal proof illustrates a broader challenge in
combining spatial constraint reasoning with description logic
blocking arguments: the global consistency properties needed for model
construction require relational information that type-based blocking
necessarily quotients away.  Resolving this tension---whether through
enriched blocking conditions, patchwork-level arguments, or entirely
new techniques---remains an important open problem.

\medskip
\noindent
The decidability of $\ALCIRCC{5}$ remains \textbf{open}, but the
triangle-type approach brings us closer to a resolution than any
previous attempt.


%% ============================================================
\section*{Acknowledgments}

This research was prompted by Michael Wessel, who introduced the
$\ALCIRCC{}$ family during his doctoral work at the University of
Hamburg.  The triangle-type approach emerged from collaborative
analysis between Wessel and Claude, informed by GPT's incisive
review~\cite{GPTReview} of the retracted paper.  The computational
verification script was developed to test the conjecture before
attempting a formal proof---an approach that proved essential for
avoiding another false claim.


%% ============================================================
\bibliographystyle{plain}
\begin{thebibliography}{10}

\bibitem{CompanionPaper}
Claude.
\newblock On the decidability of {$\mathcal{ALCI}_{\mathit{RCC5}}$} and
  {$\mathcal{ALCI}_{\mathit{RCC8}}$}: Quasimodels meet the patchwork
  property.
\newblock Discussion paper, March 2026 (revised).

\bibitem{ExtensionGapPaper}
Claude.
\newblock Closing the extension gap: Decidability of
  {$\mathcal{ALCI}_{\mathit{RCC5}}$} via direct model construction.
\newblock Discussion paper, April 2026.
\newblock \emph{Retracted}: Theorem~5.5 is false; see~\cite{GPTReview}.

\bibitem{GPTReview}
GPT.
\newblock Technical note on \emph{Closing the Extension Gap}.
\newblock Review note, April 2026.

\bibitem{Response}
Claude.
\newblock Response to GPT's review of \emph{Closing the Extension Gap}.
\newblock Discussion note, April 2026.

\bibitem{BaaderEtAl2003}
F.~Baader, D.~Calvanese, D.~McGuinness, D.~Nardi, and P.~Patel-Schneider,
  editors.
\newblock {\em The Description Logic Handbook: Theory, Implementation, and
  Applications}.
\newblock Cambridge University Press, 2003.

\bibitem{Renz1999}
J.~Renz.
\newblock Maximal tractable fragments of the {Region Connection Calculus}: A
  complete analysis.
\newblock In {\em Proc.\ IJCAI}, pages 448--454, 1999.

\bibitem{RenzNebel1999}
J.~Renz and B.~Nebel.
\newblock On the complexity of qualitative spatial reasoning: A maximal
  tractable fragment of the {Region Connection Calculus}.
\newblock {\em Artificial Intelligence}, 108(1-2):69--123, 1999.

\bibitem{Wessel2000}
M.~Wessel.
\newblock Decidable and undecidable extensions of {$\mathcal{ALC}$} with
  composition-based role inclusion axioms.
\newblock Technical Report FBI-HH-M-301/01, University of Hamburg, 2000.

\bibitem{Wessel2003}
M.~Wessel.
\newblock Qualitative spatial reasoning with the
  {$\mathcal{ALCI}_\mathit{RCC}$} family---{F}irst results and unanswered
  questions.
\newblock Technical Report FBI-HH-M-324/03, University of Hamburg, 2002/2003.

\end{thebibliography}

\end{document}
